%% annex/p07/arithmetic-proof-theory-supplement.tex
%% Editorial proof for the finite-set propositional support presentation.

\section{The finite propositional ordinal assignment}

\begin{annexitem}{annex:p7.arithmetic.1}{thm:p07-prop-cut-termination}{Termination of the finite-set propositional Cut strategy}
\begin{proof}
If \(D\) contains a nonclean Cut, the strategy chooses the first one.  Because
the cut formula occurs in a side context, one of its premises enlarges by
monotonicity to the Cut conclusion; replacing the Cut by that derivation
removes its ordinal summand and introduces no Cut.  It may also discard other
summands with the unused premise, so \(o\) strictly decreases.  Repeat until
every remaining Cut is clean.  If one remains, let \(r\) be the greatest rank
of a Cut formula and choose as specified a rank-\(r\) instance
\[
 \begin{array}{c}
 D_0:\Gamma\Rightarrow\Delta\cup\{A\}
 \qquad
 D_1:\Pi\cup\{A\}\Rightarrow\Sigma\\[-2pt]
 \hline
 \Gamma\cup\Pi\Rightarrow\Delta\cup\Sigma\quad(\mathrm{Cut}_A)
 \end{array}
 \tag{*}
\]
whose premise subtrees contain no rank-\(r\) Cut.  Such a node exists because
the tree is finite: among the rank-\(r\) Cut nodes choose one with none strictly
above it.  We verify preservation of the endsequent and decrease of the
ordinal in every possible last-rule configuration.

If one premise of (*) is a contexted identity on \(A\), erase the Cut and
retain the other premise; monotonicity supplies exactly the union contexts in
the conclusion.  If a premise is instead a contexted identity on some
\(B\ne A\), then \(B\) remains on both sides of the Cut conclusion, which is
itself a contexted identity; replace the whole Cut by that zero-premise
instance.  These are all identity leaves.  If the displayed occurrence of \(A\) was introduced only by
monotonicity in \(D_0\), discard \(D_1\), retain the monotonicity premise, and
enlarge it to the endsequent of (*); the antecedent case in \(D_1\) is
symmetric.  If a zero-premise \(\bot L\) has \(\bot\) outside the Cut
occurrence, that \(\bot\) remains in the endsequent, so \(\bot L\) closes it
directly.  When the Cut formula itself is the principal \(\bot\) of that leaf,
there is no right-introduction rule for \(\bot\); commute the Cut through the
opposite premise until one of the preceding leaf or weakening cases applies.
Each erasing case deletes the selected ordinal summand.

Suppose next that the cut occurrence is not principal in the last inference
of \(D_0\).  Apply that inference only after cutting each of its premises with
\(D_1\); if the inference has two premises, this duplicates \(D_1\).  The
support representation can fuse the removed occurrence of \(A\) with an
immediate subformula required by that other principal inference.  Before
reapplying the inference, use monotonicity to restore and align every such
displayed subformula in the affected premise.  This adds no Cut.  The result
has the same set-valued endsequent as (*).  Each new Cut
has formula rank \(r\) and a strictly smaller premise-height sum because its
left premise is a proper subtree of \(D_0\).  Any Cut nodes copied with \(D_1\)
have rank below \(r\), by the selection of (*).  The symmetric construction
applies when the cut occurrence is not principal in the last inference of
\(D_1\).  This covers every logical rule with a different principal formula,
monotonicity on another formula, and a last Cut of lower rank.

There is no contraction permutation case: repeated occurrences have one set
member and contraction is an identity of the selected context representation.
This is the point at which the support presentation avoids the invalid nested
same-rank Cut reduction for explicit multiset contraction.

It remains to consider a Cut formula principal in both last logical
inferences.  There are exactly three compound cases in the selected grammar.
For \(A=B\land C\), write the premises of \(\land R\) as
\(D_B:\Gamma\Rightarrow\Delta\cup\{B\}\) and
\(D_C:\Gamma\Rightarrow\Delta\cup\{C\}\), and the premise of \(\land L\) as
\(E:\{B,C\}\cup\Pi\Rightarrow\Sigma\).  Cut \(D_B\) with \(E\) on \(B\),
obtaining
\(\Gamma\cup\{C\}\cup\Pi\Rightarrow\Delta\cup\Sigma\).  Cut \(D_C\) with
that result on \(C\).  The repeated copies of \(\Gamma\) and \(\Delta\) are
identical set members, so the final sequent is the endsequent of (*).  Both new
cut formulas are proper subformulas of \(A\).

For \(A=B\lor C\), let
\(D:\Gamma\Rightarrow\Delta\cup\{B,C\}\) be the premise of \(\lor R\), and
let \(E_B:\{B\}\cup\Pi\Rightarrow\Sigma\) and
\(E_C:\{C\}\cup\Pi\Rightarrow\Sigma\) be the premises of \(\lor L\).
Cut \(D\) with \(E_B\) on \(B\), obtaining
\(\Gamma\cup\Pi\Rightarrow\Delta\cup\{C\}\cup\Sigma\).  Cut that result
with \(E_C\) on \(C\).  Set union collapses the repeated copies of
\(\Pi,\Sigma\), so the endsequent is unchanged and both new cut formulas have
lower rank.

For \(A=B\to C\), let
\(D:\{B\}\cup\Gamma\Rightarrow\Delta\cup\{C\}\) be the premise of
\(\to R\), and let
\(E_B:\Pi\Rightarrow\Sigma\cup\{B\}\) and
\(E_C:\{C\}\cup\Lambda\Rightarrow\Omega\) be the two premises of
\(\to L\).  First Cut \(E_B\) with \(D\) on \(B\), obtaining
\(\Pi\cup\Gamma\Rightarrow\Sigma\cup\Delta\cup\{C\}\); then Cut that
derivation with \(E_C\) on \(C\).  The result is
\(\Pi\cup\Gamma\cup\Lambda\Rightarrow\Sigma\cup\Delta\cup\Omega\), the
endsequent of (*), and both replacement Cuts have lower rank.  Atomic formulas have no
logical introduction rules, and \(\bot\) has no right-introduction rule, so
there are no further principal pairs.

We now check the ordinal rather than merely the rank of the selected Cut.
All untouched summands of \(o(D)\) remain unchanged.  A commuting or
monotonicity reduction replaces
\(\omega^{\omega r+h}\) by finitely many summands
\(\omega^{\omega r+h_i}\) with \(h_i<h\), together possibly with summands
whose exponents are below \(\omega r\).  A principal reduction replaces it
only by lower-rank summands, and an erasing reduction replaces it by none.
In every case the coefficient of the greatest exponent at which the old and
new natural sums differ drops by one.  Hence \(o(D')<o(D)\).

There is therefore no infinite sequence of selected reductions.  When the
process stops, no Cut remains: otherwise the finite tree would again contain
an uppermost Cut of greatest rank.

Finally, take any multiset derivation on \texttt{sheet:lk-defneg}.  Induction
on that tree sends every context to its support; exchange and contraction
become identities, primitive Cut becomes support Cut, every multiset Mix
becomes support Cut on the same formula, and every logical or weakening rule
becomes its set-support instance.  Nonclean Cut images are erased by
monotonicity before the reduction.
Apply the terminating reduction just proved.  A second induction chooses one
occurrence for each set member and realizes every set rule by the corresponding
multiset rule, adding only weakening and contraction when a set union merged
copies.  Because the set derivation is cut-free, this return translation is
also cut-free and has the original multiset endsequent after the required
structural adjustments.
\end{proof}
\end{annexitem}
